% ============================================================================
% Exercises for Lesson 16: Practical Projects
% Topic: LaTeX
% Solutions to capstone practice problems integrating techniques from
% all previous lessons.
% Compile: pdflatex exercises/LaTeX/16_practical_projects.tex
% ============================================================================

\documentclass[12pt, a4paper]{article}
\usepackage[utf8]{inputenc}
\usepackage[T1]{fontenc}
\usepackage{lmodern}
\usepackage[english]{babel}
\usepackage[margin=1in]{geometry}
\usepackage{microtype}

% Mathematics
\usepackage{amsmath, amssymb, amsthm}

% Figures and tables
\usepackage{graphicx}
\usepackage{float}
\usepackage{booktabs}
\usepackage{multirow}

% Lists and code
\usepackage{enumitem}
\usepackage{listings}
\usepackage{xcolor}

% Cross-references
\usepackage{cleveref}

% hyperref last (before cleveref)
\usepackage[
    colorlinks=true,
    linkcolor=blue,
    citecolor=teal,
    urlcolor=cyan
]{hyperref}

% ============================================================================
% Custom color definitions
\definecolor{codeblue}{RGB}{0,70,140}
\definecolor{codegreen}{RGB}{20,120,50}
\definecolor{codegray}{RGB}{100,100,100}

% Listings style
\lstset{
    basicstyle=\ttfamily\small,
    keywordstyle=\color{codeblue}\bfseries,
    commentstyle=\color{codegreen}\itshape,
    stringstyle=\color{red!70!black},
    numbers=left,
    numberstyle=\tiny\color{codegray},
    numbersep=8pt,
    frame=single,
    breaklines=true,
    showstringspaces=false,
    tabsize=4
}

% Theorem environment
\newtheorem{theorem}{Theorem}[section]
\newtheorem{lemma}[theorem]{Lemma}

% ============================================================================

\title{Exercises for Lesson 16: Practical Projects}
\author{Study Project}
\date{}

\begin{document}
\maketitle

\tableofcontents
\bigskip

\noindent
This exercise file integrates skills from all previous lessons into four
complete project tasks.  Each task asks you to build a realistic, compilable
\LaTeX{} document that combines multiple packages and techniques.

% ============================================================================
% PROJECT 1: ACADEMIC PAPER
% Problem: A complete academic paper with math, figures, tables, bibliography,
%          and cross-references.
% ============================================================================

\section*{Project 1: Complete Academic Paper}
\addcontentsline{toc}{section}{Project 1: Complete Academic Paper}

\textbf{Goal:} Produce a 4--6 page conference-style paper that contains at
least one section with numbered equations and cross-references, one
floating figure with a meaningful caption, one booktabs table, and a
bibliography with at least three citations (BibLaTeX or BibTeX).

\subsection*{Solution}

The solution demonstrates a condensed version.  Study each region before
building your own full version.

\begin{verbatim}
% ============================================================
% academic_paper.tex
% Compile: pdflatex -> biber -> pdflatex -> pdflatex
% ============================================================
\documentclass[11pt,a4paper]{article}

\usepackage[utf8]{inputenc}
\usepackage[T1]{fontenc}
\usepackage{lmodern}
\usepackage[english]{babel}
\usepackage[margin=1in]{geometry}
\usepackage{microtype}
\usepackage{amsmath, amssymb, amsthm}
\usepackage{graphicx}
\usepackage{booktabs}
\usepackage[backend=biber,style=authoryear]{biblatex}
\addbibresource{paper_refs.bib}
\usepackage{hyperref}
\usepackage{cleveref}

% Theorem environments
\newtheorem{theorem}{Theorem}[section]
\newtheorem{proposition}[theorem]{Proposition}

% Custom commands
\newcommand{\R}{\mathbb{R}}
\newcommand{\norm}[1]{\lVert#1\rVert}

\title{Gradient Descent Convergence in Strongly Convex Landscapes}
\author{Jane Doe \and John Smith}
\date{\today}

\begin{document}
\maketitle

\begin{abstract}
We study the convergence rate of gradient descent (GD) applied to
strongly convex, smooth objective functions.  We derive a tight
exponential decay bound and verify it empirically on a least-squares
regression benchmark.  Our results confirm the $O(\rho^k)$ convergence
with $\rho = 1 - \mu/L$, where $\mu$ and $L$ are the strong convexity
and Lipschitz constants, respectively.
\end{abstract}

\section{Introduction}
\label{sec:intro}

Gradient descent is the workhorse of modern machine learning
\parencite{goodfellow2016}.  For a differentiable function
$f: \R^n \to \R$, the GD update at iteration $k$ is
\begin{equation}
    \label{eq:gd_update}
    \mathbf{x}_{k+1} = \mathbf{x}_k - \alpha \nabla f(\mathbf{x}_k),
\end{equation}
where $\alpha > 0$ is the step size.  \Cref{sec:theory} establishes our
convergence theorem; \cref{sec:experiments} confirms it numerically.

\section{Theoretical Analysis}
\label{sec:theory}

\begin{theorem}[Convergence of GD on Strongly Convex Functions]
\label{thm:gd_conv}
Let $f: \R^n \to \R$ be $\mu$-strongly convex and $L$-smooth with
minimiser $\mathbf{x}^*$.  With step size $\alpha = 1/L$, the GD iterates
satisfy
\begin{equation}
    \label{eq:conv_bound}
    \norm{\mathbf{x}_k - \mathbf{x}^*}^2
    \leq \left(1 - \frac{\mu}{L}\right)^k
         \norm{\mathbf{x}_0 - \mathbf{x}^*}^2.
\end{equation}
\end{theorem}

\begin{proof}
By $L$-smoothness and the update rule \eqref{eq:gd_update},
\[
    f(\mathbf{x}_{k+1}) \leq f(\mathbf{x}_k)
    - \frac{1}{2L}\norm{\nabla f(\mathbf{x}_k)}^2.
\]
Combining with the $\mu$-strong convexity inequality
$f(\mathbf{x}_k) - f(\mathbf{x}^*) \geq (\mu/2)\norm{\mathbf{x}_k - \mathbf{x}^*}^2$
and iterating completes the proof \parencite[Ch.\,3]{boyd2004}.
\end{proof}

\section{Numerical Experiments}
\label{sec:experiments}

We verify \cref{thm:gd_conv} on least-squares regression with a random
SPD system.  \cref{fig:convergence} shows the convergence curve; the
dashed line is the theoretical bound from \cref{eq:conv_bound}.

\begin{figure}[h]
\centering
% A placeholder box; replace with \includegraphics{convergence.pdf}
\fbox{\parbox{0.55\textwidth}{\centering\vspace{2cm}
      \textit{[Convergence curve figure]}\vspace{2cm}}}
\caption{GD convergence on a $500 \times 500$ random SPD system.
         Solid = empirical; dashed = bound from \cref{eq:conv_bound}.}
\label{fig:convergence}
\end{figure}

\cref{tab:results} summarises results for three problem sizes.

\begin{table}[h]
\centering
\caption{Iterations to convergence ($\norm{\mathbf{x}_k - \mathbf{x}^*}
         / \norm{\mathbf{x}_0 - \mathbf{x}^*} \leq 10^{-6}$).}
\label{tab:results}
\begin{tabular}{lrrrr}
\toprule
\textbf{Size} & \boldmath$\mu$ & \boldmath$L$ & \boldmath$\kappa = L/\mu$
    & \textbf{Iterations} \\
\midrule
$n = 100$  & 0.52 & 12.3 & 23.7  & 312  \\
$n = 500$  & 0.11 & 18.7 & 170.0 & 1840 \\
$n = 1000$ & 0.04 & 21.5 & 537.5 & 5819 \\
\bottomrule
\end{tabular}
\end{table}

\section{Conclusion}
\label{sec:conclusion}

We have verified the classical $O(\rho^k)$ convergence of gradient descent
on strongly convex functions (\cref{thm:gd_conv}).  Future work will
extend the analysis to stochastic and momentum-based variants
\parencite{ruder2016}.

\printbibliography

\end{document}
\end{verbatim}

\medskip
\noindent
\textbf{Key techniques demonstrated:}
\begin{itemize}[noitemsep]
    \item BibLaTeX (\texttt{authoryear}) with \verb|\parencite| and
          \verb|\printbibliography|
    \item Numbered equation (\texttt{equation}) with \verb|\label|/\verb|\eqref|
    \item Theorem environment with a complete proof
    \item Floating figure and booktabs table, both cross-referenced with
          \verb|\cref|
    \item Custom commands (\verb|\R|, \verb|\norm|)
    \item \texttt{cleveref} \verb|\Cref| / \verb|\cref| for all references
\end{itemize}

% ============================================================================
% PROJECT 2: BEAMER PRESENTATION WITH TIKZ AND OVERLAYS
% Problem: A Beamer slide deck with TikZ diagrams and overlays.
% ============================================================================

\section*{Project 2: Beamer Presentation with TikZ and Overlays}
\addcontentsline{toc}{section}{Project 2: Beamer Presentation}

\textbf{Goal:} Produce a 6--8 slide presentation on a technical topic that
includes: a title slide, a table of contents slide, at least two slides with
\textit{overlay specifications} (\verb|\pause|, \verb|\only|,
\verb|\uncover|), at least one slide with a TikZ diagram, and a summary slide.

\subsection*{Solution}

\begin{verbatim}
% ============================================================
% presentation.tex
% Compile: pdflatex presentation.tex (twice)
% ============================================================
\documentclass[12pt]{beamer}

% Theme and colour
\usetheme{Madrid}
\usecolortheme{seahorse}

% Packages
\usepackage[utf8]{inputenc}
\usepackage[T1]{fontenc}
\usepackage{amsmath, amssymb}
\usepackage{tikz}
\usetikzlibrary{arrows.meta, shapes.geometric, positioning}

% Metadata
\title{Gradient Descent for Machine Learning}
\subtitle{Theory, Variants, and Practice}
\author{Jane Doe}
\institute{Department of Computer Science\\Example University}
\date{\today}

\begin{document}

% ----- Slide 1: Title -----
\begin{frame}
\titlepage
\end{frame}

% ----- Slide 2: Table of Contents -----
\begin{frame}{Outline}
\tableofcontents
\end{frame}

% ----- Section 1 -----
\section{Background}

% ----- Slide 3: Background with overlays -----
\begin{frame}{Optimisation Background}
We want to minimise $f: \mathbb{R}^n \to \mathbb{R}$.

\medskip
\begin{itemize}
    \item<1-> \textbf{Objective:} find $\mathbf{x}^* = \arg\min f(\mathbf{x})$
    \item<2-> \textbf{Gradient:} $\nabla f(\mathbf{x})$ points uphill
    \item<3-> \textbf{Insight:} move in direction $-\nabla f(\mathbf{x})$
    \item<4-> \textbf{Update rule:}
        $\mathbf{x}_{k+1} = \mathbf{x}_k - \alpha \nabla f(\mathbf{x}_k)$
\end{itemize}

\onslide<5->{
\begin{block}{Key parameter}
Step size $\alpha$ controls the trade-off between speed and stability.
\end{block}
}
\end{frame}

% ----- Section 2 -----
\section{TikZ Diagram}

% ----- Slide 4: TikZ gradient descent illustration -----
\begin{frame}{Visualising Gradient Descent}
\begin{center}
\begin{tikzpicture}[scale=1.1,
    point/.style={circle, fill, inner sep=2pt},
    arrow/.style={-{Stealth}, thick, blue!70}
]
    % Contour ellipses representing a bowl-shaped function
    \draw[gray!40, thick] (0,0) ellipse (2.5cm and 1.5cm);
    \draw[gray!60, thick] (0,0) ellipse (1.8cm and 1.0cm);
    \draw[gray!80, thick] (0,0) ellipse (1.0cm and 0.6cm);

    % Optimum
    \node[point, red, label=below right:{\small $\mathbf{x}^*$}] at (0,0) {};

    % Gradient descent path
    \node[point, blue, label=above:{\small $\mathbf{x}_0$}] (x0) at (2.2,1.2) {};
    \node[point, blue, label=above:{\small $\mathbf{x}_1$}] (x1) at (1.4,0.6) {};
    \node[point, blue, label=right:{\small $\mathbf{x}_2$}] (x2) at (0.8,0.2) {};
    \node[point, blue] (x3) at (0.3,0.05) {};

    \draw[arrow] (x0) -- (x1);
    \draw[arrow] (x1) -- (x2);
    \draw[arrow] (x2) -- (x3);
    \draw[arrow] (x3) -- (0.05,0.01);
\end{tikzpicture}
\end{center}
\vspace{-0.5em}
Each arrow $\mathbf{x}_k \to \mathbf{x}_{k+1}$ is proportional to
$-\alpha\,\nabla f(\mathbf{x}_k)$.
\end{frame}

% ----- Section 3 -----
\section{Variants}

% ----- Slide 5: Variants with \only -----
\begin{frame}{GD Variants}
\only<1>{
\textbf{Batch Gradient Descent (GD)}
\begin{itemize}
    \item Uses all training examples per update
    \item Stable convergence
    \item Slow on large datasets ($O(n)$ per step)
\end{itemize}
}
\only<2>{
\textbf{Stochastic Gradient Descent (SGD)}
\begin{itemize}
    \item Uses a single example per update
    \item Fast; can escape local minima
    \item Noisy gradients; needs learning rate schedule
\end{itemize}
}
\only<3>{
\textbf{Mini-Batch SGD}
\begin{itemize}
    \item Uses a small batch (32--512 examples) per update
    \item Best of both: GPU-friendly, reasonable variance
    \item Standard in deep learning practice
\end{itemize}
}
\end{frame}

% ----- Slide 6: Summary with \uncover -----
\begin{frame}{Summary}
\begin{enumerate}
    \item<1-> Gradient descent minimises $f$ by following $-\nabla f$
    \item<2-> Convergence rate is $O\!\left(\left(1-\frac{\mu}{L}\right)^k\right)$
              for strongly convex $f$
    \item<3-> Three variants: GD, SGD, mini-batch SGD
    \item<4-> Choice of step size $\alpha$ is critical

    \uncover<5->{
    \item[\checkmark] Next: \textbf{momentum methods} (Nesterov, Adam)
    }
\end{enumerate}
\end{frame}

\end{document}
\end{verbatim}

\medskip
\noindent
\textbf{Key techniques demonstrated:}
\begin{itemize}[noitemsep]
    \item \texttt{beamer} theme/colour scheme selection
    \item Overlay specifications: \verb|<1->|, \verb|\only<n>|,
          \verb|\uncover<n->|, \verb|\pause|
    \item TikZ diagram embedded in a \texttt{frame}
    \item \texttt{block} environment for highlighted boxes
    \item \verb|\tableofcontents| on a dedicated slide
\end{itemize}

% ============================================================================
% PROJECT 3: MULTI-FILE REPORT STRUCTURE
% Problem: A multi-file report with \input / \include.
% ============================================================================

\section*{Project 3: Multi-File Report Project}
\addcontentsline{toc}{section}{Project 3: Multi-File Report Structure}

\textbf{Goal:} Organise a report as a \emph{multi-file project}: one main
file that loads shared preamble settings and individual chapter files via
\verb|\include|.  The structure should mirror a real-world report where each
author edits their chapter independently.

\subsection*{Recommended Directory Layout}

\begin{verbatim}
report/
├── main.tex            % root file: loads preamble + chapters
├── preamble.tex        % shared package list and settings
├── chapters/
│   ├── 01_intro.tex
│   ├── 02_methods.tex
│   ├── 03_results.tex
│   └── 04_conclusion.tex
├── figures/
│   └── diagram.pdf
└── references.bib
\end{verbatim}

\subsection*{main.tex}

\begin{verbatim}
% main.tex — root file
\documentclass[12pt,a4paper]{report}
\input{preamble}           % shared settings

\title{Annual Research Report 2024}
\author{Research Team Alpha}
\date{\today}

\begin{document}

\maketitle

\pagenumbering{roman}
\tableofcontents
\listoffigures
\listoftables
\clearpage
\pagenumbering{arabic}

\include{chapters/01_intro}
\include{chapters/02_methods}
\include{chapters/03_results}
\include{chapters/04_conclusion}

\bibliographystyle{plain}
\bibliography{references}

\end{document}
\end{verbatim}

\subsection*{preamble.tex}

\begin{verbatim}
% preamble.tex — shared package declarations
\usepackage[utf8]{inputenc}
\usepackage[T1]{fontenc}
\usepackage{lmodern}
\usepackage[english]{babel}
\usepackage[margin=2.5cm]{geometry}
\usepackage{microtype}
\usepackage{amsmath, amssymb}
\usepackage{graphicx}
\usepackage{booktabs}
\usepackage{setspace}
\onehalfspacing
\usepackage{fancyhdr}
\pagestyle{fancy}
\fancyhf{}
\fancyhead[LE,RO]{\slshape\nouppercase{\leftmark}}
\fancyhead[RE,LO]{\slshape Annual Research Report}
\fancyfoot[C]{\thepage}
\usepackage{hyperref}
\usepackage{cleveref}
\end{verbatim}

\subsection*{chapters/01\_intro.tex (representative chapter)}

\begin{verbatim}
% chapters/01_intro.tex
\chapter{Introduction}
\label{chap:intro}

This report summarises the research activities of Team Alpha during 2024.
\cref{chap:methods} describes the experimental setup; \cref{chap:results}
presents the main findings.

\section{Background}
\label{sec:background}
Prior work in this area has established ...
\end{verbatim}

\medskip
\noindent
\textbf{Key techniques demonstrated:}
\begin{itemize}[noitemsep]
    \item \verb|\input| for shared preamble (processed inline)
    \item \verb|\include| for chapters (each gets its own \texttt{.aux} file,
          enabling partial compilation with \verb|\includeonly{chap1}|)
    \item Separate \texttt{preamble.tex} — change once, applies everywhere
    \item Roman + arabic page numbering switch around front matter
    \item \texttt{fancyhdr} configured centrally for the whole report
\end{itemize}

\subsection*{Partial Compilation}

To compile only Chapter 2 while still resolving all cross-references from
the full run:

\begin{verbatim}
% In main.tex, after \input{preamble}:
\includeonly{chapters/02_methods}
\end{verbatim}

\LaTeX{} reads the cached \texttt{.aux} files of the omitted chapters, so
cross-references and page numbers remain correct.

% ============================================================================
% PROJECT 4: CV WITH CUSTOM COMMANDS
% Problem: A one-page CV using custom commands for consistent formatting.
% ============================================================================

\section*{Project 4: Curriculum Vitae with Custom Commands}
\addcontentsline{toc}{section}{Project 4: CV with Custom Commands}

\textbf{Goal:} Build a clean one-page CV \emph{without} the \texttt{moderncv}
package.  Define custom \LaTeX{} commands that encapsulate the formatting so
that adding a new entry requires only a single macro call.

\subsection*{Solution}

\begin{verbatim}
% ============================================================
% cv.tex — custom CV without moderncv
% Compile: pdflatex cv.tex
% ============================================================
\documentclass[10pt,a4paper]{article}

\usepackage[utf8]{inputenc}
\usepackage[T1]{fontenc}
\usepackage{lmodern}
\usepackage[
    top=1.5cm, bottom=1.5cm,
    left=2cm,  right=2cm
]{geometry}
\usepackage{xcolor}
\usepackage{enumitem}
\usepackage{microtype}
\usepackage{hyperref}
\hypersetup{hidelinks}

% ---------- colour palette ----------
\definecolor{cvblue}{RGB}{0,56,101}
\definecolor{cvgray}{RGB}{90,90,90}

% ---------- section heading ----------
% \cvsection{Title}
\newcommand{\cvsection}[1]{%
    \vspace{0.5em}%
    {\color{cvblue}\large\bfseries #1}%
    \par\nobreak
    {\color{cvblue}\rule{\textwidth}{0.6pt}}%
    \vspace{0.3em}%
}

% ---------- entry: date, title, organisation, location ----------
% \cventry{2022--2024}{M.Sc. Computer Science}{MIT}{Cambridge, MA}
\newcommand{\cventry}[4]{%
    \noindent
    \begin{tabular*}{\textwidth}{@{}l@{\extracolsep{\fill}}r@{}}
        \textbf{#2}, #3 & {\color{cvgray}#1, #4}
    \end{tabular*}%
    \par
}

% ---------- skill row ----------
% \cvskill{Category}{Value}
\newcommand{\cvskill}[2]{%
    \noindent
    \textbf{#1:} #2\par\smallskip
}

% ---------- no paragraph indent ----------
\setlength{\parindent}{0pt}
\setlength{\parskip}{0pt}

% ============================================================
\begin{document}

% Name and contact
\begin{center}
    {\LARGE\bfseries\color{cvblue} Jane Doe}\\[4pt]
    {\color{cvgray}
     \href{mailto:jane@example.com}{jane@example.com} \quad
     \textbar\quad +1-555-0123 \quad
     \textbar\quad
     \href{https://github.com/janedoe}{github.com/janedoe}
    }
\end{center}

\bigskip

% --- Education ---
\cvsection{Education}

\cventry{2022--2024}{M.Sc.\ Computer Science}{MIT}{Cambridge, MA}
\begin{itemize}[noitemsep, topsep=2pt, leftmargin=1.5em]
    \item Thesis: \textit{Convergence of Gradient Methods for Sparse
          Optimisation}
    \item GPA: 4.0/4.0
\end{itemize}

\medskip
\cventry{2018--2022}{B.Sc.\ Mathematics (Hons.)}{Stanford}{Stanford, CA}
\begin{itemize}[noitemsep, topsep=2pt, leftmargin=1.5em]
    \item Minor in Computer Science
    \item Dean's List all semesters
\end{itemize}

% --- Experience ---
\cvsection{Work Experience}

\cventry{Jun 2024--Present}{Machine Learning Researcher}{DeepMind}{London, UK}
\begin{itemize}[noitemsep, topsep=2pt, leftmargin=1.5em]
    \item Developed novel preconditioning strategies, reducing training
          time by 30\%
    \item Co-authored two NeurIPS papers on optimisation for transformers
\end{itemize}

\medskip
\cventry{May--Aug 2023}{Research Intern}{Google Research}{Mountain View, CA}
\begin{itemize}[noitemsep, topsep=2pt, leftmargin=1.5em]
    \item Implemented efficient sparse attention kernels in CUDA
    \item Achieved 2$\times$ inference speedup on long-context benchmarks
\end{itemize}

% --- Skills ---
\cvsection{Technical Skills}

\cvskill{Languages}{Python, C++, Rust, Julia, Bash}
\cvskill{Frameworks}{PyTorch, JAX, TensorFlow, NumPy, SciPy}
\cvskill{Tools}{Git, Docker, Kubernetes, \LaTeX, Linux}

% --- Publications ---
\cvsection{Selected Publications}

\begin{enumerate}[noitemsep, topsep=2pt, leftmargin=1.5em]
    \item J.\ Doe and J.\ Smith. ``Hierarchical Preconditioning for
          Large-Scale Optimisation.'' \textit{NeurIPS 2024}.
    \item J.\ Doe et al. ``Sparse Attention for Long Documents.''
          \textit{ICML 2023}.
\end{enumerate}

\end{document}
\end{verbatim}

\medskip
\noindent
\textbf{Key techniques demonstrated:}
\begin{itemize}[noitemsep]
    \item Custom commands with arguments: \verb|\cvsection{#1}|,
          \verb|\cventry{#1}{#2}{#3}{#4}|, \verb|\cvskill{#1}{#2}|
    \item \texttt{tabular*} with \verb|\extracolsep{\fill}| for flush-right
          date/location
    \item \verb|\definecolor| and \verb|\color| for a consistent brand palette
    \item \texttt{enumitem} options (\texttt{noitemsep}, \texttt{topsep}) to
          tighten list spacing in a dense one-page layout
    \item \verb|\hyperref| with \verb|hidelinks| for clickable email/GitHub
          without visible coloured boxes
\end{itemize}

\subsection*{Extending the CV Template}

Adding a new work entry requires only:

\begin{verbatim}
\cventry{Start--End}{Role}{Company}{Location}
\begin{itemize}[noitemsep, topsep=2pt, leftmargin=1.5em]
    \item Accomplishment one
    \item Accomplishment two
\end{itemize}
\end{verbatim}

No formatting knowledge is needed once the macros are defined---the caller
only supplies content.

% ============================================================================
% CHECKLIST
% ============================================================================

\bigskip
\noindent\rule{\textwidth}{0.4pt}

\subsection*{Capstone Checklist}

Before considering each project complete, verify:

\begin{enumerate}
    \item \textbf{Project 1 (Academic Paper)}
    \begin{itemize}[noitemsep]
        \item[$\square$] Compiles with \texttt{pdflatex} $\to$ \texttt{biber} $\to$
              \texttt{pdflatex} $\to$ \texttt{pdflatex}
        \item[$\square$] At least one numbered equation with \verb|\label|/\verb|\eqref|
        \item[$\square$] At least one floating figure and one booktabs table
        \item[$\square$] At least three bibliography entries cited in text
        \item[$\square$] All figures, tables, equations, and sections
              cross-referenced with \verb|\cref|
    \end{itemize}

    \item \textbf{Project 2 (Beamer)}
    \begin{itemize}[noitemsep]
        \item[$\square$] Title slide + TOC slide
        \item[$\square$] At least two slides with overlay specifications
        \item[$\square$] At least one slide containing a TikZ diagram
        \item[$\square$] Summary/conclusion slide
    \end{itemize}

    \item \textbf{Project 3 (Multi-File Report)}
    \begin{itemize}[noitemsep]
        \item[$\square$] Separate \texttt{preamble.tex}, at least three chapter files
        \item[$\square$] Front matter with roman page numbers; main matter arabic
        \item[$\square$] Cross-chapter \verb|\cref| references resolve correctly
        \item[$\square$] Partial compilation with \verb|\includeonly| tested
    \end{itemize}

    \item \textbf{Project 4 (CV)}
    \begin{itemize}[noitemsep]
        \item[$\square$] At least three custom commands (\verb|\cvsection|,
              \verb|\cventry|, \verb|\cvskill| or equivalents)
        \item[$\square$] Consistent colour palette via \verb|\definecolor|
        \item[$\square$] Fits on one page
        \item[$\square$] All hyperlinks functional (\texttt{hyperref})
    \end{itemize}
\end{enumerate}

\end{document}
